% =========================================================
% Limits of Functions — Notes
% Chapter: functions
% Section: limits
% Volume III · Analysis
% =========================================================

\section{Limits of Functions}
\label{sec:limits-of-functions}

% ---------------------------------------------------------
% Toolkit
% ---------------------------------------------------------

\begin{tcolorbox}[title={Toolkit --- Limits of Functions}]
\begin{itemize}
  \item $\lim_{x \to c} f(x) = L$ requires $c$ to be a cluster point
    of $\operatorname{dom}(f)$. The value $f(c)$ is irrelevant;
    $f$ need not be defined at $c$.
  \item The input condition is punctured:
    $\operatorname{InputTolerance}(x,c,\delta)$ means
    $0 < |x-c| < \delta$. This is what separates limits from
    continuity, where $\operatorname{Within}(x,c,\delta)$
    permits $x = c$.
  \item Three equivalent forms: $\varepsilon$-$\delta$, neighbourhood,
    sequential. The Sequential Criterion is the primary working tool.
  \item Bridge principle: algebra of sequence limits transfers
    automatically via the Sequential Criterion.
  \item A finite limit implies $\operatorname{BoundedNear}(f,c)$.
  \item Limits are local: behavior far from $c$ is irrelevant.
\end{itemize}
\end{tcolorbox}

% ---------------------------------------------------------
% Tao scaffolding
% ---------------------------------------------------------

\begin{definition}[$\varepsilon$-Closeness of a Function to a Value]
\label{def:eps-close-function}
Let $X \subseteq \mathbb{R}$, let $f : X \to \mathbb{R}$, let
$L \in \mathbb{R}$, and let $\varepsilon > 0$. The function $f$ is
\textbf{$\varepsilon$-close to $L$} if
$\operatorname{OutputTolerance}(f(x), L, \varepsilon)$ holds for all
$x \in X$, i.e., $|f(x) - L| \leq \varepsilon$ for all $x \in X$.
\end{definition}

\begin{remark*}[Standard quantified statement]
\[
  f \;\varepsilon\text{-close to } L
  \;\iff\;
  \forall x \in X :\;
  \operatorname{OutputTolerance}(f(x), L, \varepsilon).
\]
\end{remark*}

\begin{remark*}[Interpretation]
A global condition: every output lies within $\varepsilon$ of $L$,
regardless of where in $X$ the input lies. This is Tao's scaffolding
toward the limit definition, which localises the same output
condition to a punctured neighbourhood of $c$.
\end{remark*}

\begin{definition}[Local $\delta$-Closeness Near a Point]
\label{def:delta-close-function}
Let $f : X \to \mathbb{R}$, let $L \in \mathbb{R}$, let $x_0$ be
adherent to $X$, and let $\varepsilon, \delta > 0$. The function $f$
is \textbf{$\varepsilon$-close to $L$ near $x_0$ within $\delta$} if
\[
  \forall x \in X :\;
  \operatorname{Within}(x, x_0, \delta)
  \;\Rightarrow\;
  \operatorname{OutputTolerance}(f(x), L, \varepsilon).
\]
\end{definition}

\begin{remark*}[Standard quantified statement]
\[
  \forall x \in X :\;
  |x - x_0| < \delta \;\Rightarrow\; |f(x) - L| \leq \varepsilon.
\]
\end{remark*}

\begin{remark*}[Interpretation]
A local condition: only outputs from the $\delta$-ball around $x_0$
need lie within $\varepsilon$ of $L$. The input condition here is
$\operatorname{Within}$ — unpunctured. The limit definition replaces
this with $\operatorname{InputTolerance}$ (punctured) and asserts
that for every $\varepsilon$ such a $\delta$ exists.
\end{remark*}

% ---------------------------------------------------------
% Definition — Limit of a Function
% ---------------------------------------------------------

\begin{tcolorbox}[title={Definition --- Limit of a Function}]
\begin{definition}[Limit of a Function]
\label{def:limit-function}
Let $f : A \to \mathbb{R}$ and let $c$ be a cluster point of $A$.
A real number $L$ is the \textbf{limit of $f$ at $c$}, written
$\lim_{x \to c} f(x) = L$, if
\[
  \ForallOT{\varepsilon},\;
  \ExistsIT{\delta},\;
  \forall x \in A :\;
  \operatorname{InputTolerance}(x, c, \delta)
  \;\Rightarrow\;
  \operatorname{OutputTolerance}(f(x), L, \varepsilon).
\]
If no such $L$ exists, $f$ \textbf{diverges} at $c$.

\smallskip\noindent
\hyperref[prf:limit-unique]{\textit{Go to uniqueness proof.}}
\end{definition}
\end{tcolorbox}

\begin{remark*}[Standard quantified statement]
\[
  \operatorname{FunctionLimit}(f, c, L)
  \;\iff\;
  \forall\varepsilon > 0,\;
  \exists\delta > 0,\;
  \forall x \in A :\;
  0 < |x - c| < \delta
  \;\Rightarrow\;
  |f(x) - L| < \varepsilon.
\]
\end{remark*}

\begin{remark*}[Definition predicate reading]
\[
  \operatorname{FunctionLimit}(f, c, L)
  \;\colonequiv\;
  \ForallOT{\varepsilon},\;
  \ExistsIT{\delta},\;
  \forall x \in A :\;
  \operatorname{InputTolerance}(x, c, \delta)
  \;\Rightarrow\;
  \operatorname{OutputTolerance}(f(x), L, \varepsilon).
\]
\end{remark*}

\begin{remark*}[Negated quantified statement]
\[
  \neg\,\operatorname{FunctionLimit}(f, c, L)
  \;\iff\;
  \exists\varepsilon_0 > 0,\;
  \forall\delta > 0,\;
  \exists x \in A :\;
  \operatorname{InputTolerance}(x, c, \delta)
  \;\wedge\;
  \neg\,\operatorname{OutputTolerance}(f(x), L, \varepsilon_0).
\]
\end{remark*}

\begin{remark*}[Negation predicate reading]
\[
  \neg\,\operatorname{FunctionLimit}(f, c, L)
  \;\colonequiv\;
  \exists\varepsilon_0 > 0,\;
  \forall\delta > 0,\;
  \exists x \in A :\;
  0 < |x - c| < \delta
  \;\wedge\;
  |f(x) - L| \geq \varepsilon_0.
\]
\end{remark*}

\begin{remark*}[Failure mode decomposition]
The negation says: some fixed output gap $\varepsilon_0 > 0$ cannot
be closed no matter how small the input window.
\[
  \underbrace{%
    \exists\,\text{two sequences } x_n, y_n \to c
    \text{ with } f(x_n) \to L_1 \neq L_2 \leftarrow f(y_n)
  }_{\text{limit does not exist at all}}
  \quad\vee\quad
  \underbrace{%
    f(x_n) \text{ oscillates for every } x_n \to c
  }_{\text{essential oscillation near } c}.
\]
\end{remark*}

\begin{remark*}[Interpretation]
The quantifier structure is asymmetric by design. The output
tolerance $\varepsilon$ is a \emph{demand}: the caller specifies how
precise the output must be. The input tolerance $\delta$ is a
\emph{response}: the definition asserts one always exists. The input
condition is punctured ($\operatorname{InputTolerance}$, not
$\operatorname{Within}$) because $f$ may not be defined at $c$, and
because continuity — not limits — is the concept that includes $c$
itself. This is the deepest structural difference between the two
definitions. Four critical features: (i) $f$ need not be defined at
$c$; (ii) $c$ must be a cluster point of $A$ else the condition is
vacuously satisfied by every $L$; (iii) $\delta$ may depend on
$\varepsilon$; (iv) the condition is about all $x$ near $c$
simultaneously, not just one.
\end{remark*}

% ---------------------------------------------------------
% Definition — Sequential Form of Limit
% ---------------------------------------------------------

\begin{definition}[Sequential Definition of Limit]
\label{def:limit-function-sequential}
Let $f : A \to \mathbb{R}$ and let $c$ be a cluster point of $A$.
The limit $\lim_{x \to c} f(x) = L$ in the \textbf{sequential
sense} holds if for every sequence $(x_n) \subseteq A \setminus \{c\}$
with $x_n \to c$, the sequence $(f(x_n))$ converges to $L$.
\end{definition}

\begin{remark*}[Standard quantified statement]
\[
  \operatorname{FunctionLimit}(f, c, L) \text{ (sequential)}
  \;\iff\;
  \forall (x_n) \subseteq A \setminus \{c\} :\;
  x_n \to c \;\Rightarrow\; f(x_n) \to L.
\]
\end{remark*}

\begin{remark*}[Definition predicate reading]
\[
  \operatorname{FunctionLimit}_{\mathrm{seq}}(f, c, L)
  \;\colonequiv\;
  \forall (x_n) \subseteq A \setminus \{c\} :\;
  \operatorname{ConvergentSequence}(x_n, c)
  \;\Rightarrow\;
  \operatorname{ConvergentSequence}(f(x_n), L).
\]
\end{remark*}

\begin{remark*}[Negated quantified statement]
\[
  \neg\,\operatorname{FunctionLimit}_{\mathrm{seq}}(f, c, L)
  \;\iff\;
  \exists (x_n) \subseteq A \setminus \{c\} :\;
  x_n \to c \;\wedge\; f(x_n) \not\to L.
\]
\end{remark*}

\begin{remark*}[Interpretation]
The sequence $(x_n)$ must lie in $A \setminus \{c\}$: the term
$x_n = c$ is excluded, mirroring the $0 <$ in the $\varepsilon$-$\delta$
form. The sequential form is the primary tool for both proving limits
(via the bridge principle) and disproving them (exhibit one bad
sequence). The $\varepsilon$-$\delta$ form is better for constructing
explicit $\delta$ functions.
\end{remark*}

% ---------------------------------------------------------
% Proposition — Neighbourhood Form
% ---------------------------------------------------------

\begin{proposition}[Neighbourhood Form of the Limit]
\label{prop:limit-neighbourhood-equiv}
Let $f : A \to \mathbb{R}$ and let $c$ be a cluster point of $A$.
The following are equivalent:
\begin{enumerate}[label=(\roman*)]
  \item $\lim_{x \to c} f(x) = L$.
  \item For every $\varepsilon$-neighbourhood $V_\varepsilon(L)$
    there exists a $\delta$-neighbourhood $V_\delta(c)$ such that
    $x \in V_\delta^*(c) \cap A \Rightarrow f(x) \in V_\varepsilon(L)$.
\end{enumerate}

\smallskip\noindent
\hyperref[prf:limit-neighbourhood-equiv]{\textit{Go to proof.}}
\end{proposition}

\begin{remark*}[Standard quantified statement]
\[
\begin{aligned}
(i)\quad &\ForallOT{\varepsilon},\;
  \ExistsIT{\delta},\;
  \forall x \in A :\;
  \operatorname{InputTolerance}(x,c,\delta)
  \Rightarrow
  \operatorname{OutputTolerance}(f(x),L,\varepsilon). \\[4pt]
(ii)\quad &\forall\varepsilon>0,\;
  \exists\delta>0,\;
  \forall x \in V_\delta^*(c) \cap A :\;
  f(x) \in V_\varepsilon(L).
\end{aligned}
\]
\end{remark*}

\begin{remark*}[Interpretation]
The equivalence is purely notational: $V_\delta^*(c) \cap A$ is
exactly the set of $x \in A$ satisfying
$\operatorname{InputTolerance}(x,c,\delta)$, and $f(x) \in
V_\varepsilon(L)$ is exactly $\operatorname{OutputTolerance}(f(x),L,
\varepsilon)$. The neighbourhood form makes the geometric content
explicit: the image of every punctured $\delta$-slice of $A$ must
fit inside the $\varepsilon$-ball around $L$. This is the form of
continuity used in metric space topology.
\end{remark*}

% ---------------------------------------------------------
% Theorem — Uniqueness of Limits
% ---------------------------------------------------------

\begin{theorem}[Uniqueness of Limits of Functions]
\label{thm:limit-unique}
Let $f : A \to \mathbb{R}$ and let $c$ be a cluster point of $A$.
Then $f$ has at most one limit at $c$.

\smallskip\noindent
\hyperref[prf:limit-unique]{\textit{Go to proof.}}
\end{theorem}

\begin{remark*}[Standard quantified statement]
\[
  \operatorname{FunctionLimit}(f,c,L)
  \;\wedge\;
  \operatorname{FunctionLimit}(f,c,L')
  \;\Rightarrow\; L = L'.
\]
\end{remark*}

\begin{remark*}[Failure mode decomposition]
Non-uniqueness would require two values $L \neq L'$ each satisfying
the limit definition. The $\varepsilon = |L-L'|/2$ argument shows
any $x$ within $\min(\delta_1,\delta_2)$ of $c$ satisfies both
$|f(x)-L| < \varepsilon$ and $|f(x)-L'| < \varepsilon$, but then
$|L-L'| \leq |L-f(x)|+|f(x)-L'| < 2\varepsilon = |L-L'|$, a
contradiction. The cluster point condition is essential: it guarantees
such an $x$ exists.
\end{remark*}

\begin{remark*}[Interpretation]
The limit, when it exists, is uniquely determined by the local
behavior of $f$ near $c$. Uniqueness is what licenses writing
$\lim_{x\to c} f(x) = L$ as an equation rather than an assertion.
\end{remark*}

% ---------------------------------------------------------
% Proposition — Sequential Criterion
% ---------------------------------------------------------

\begin{proposition}[Sequential Criterion for Function Limits]
\label{prop:sequential-criterion-limits}
Let $f : A \to \mathbb{R}$ and let $c$ be a cluster point of $A$.
The following are equivalent:
\begin{enumerate}[label=(\alph*)]
  \item $\lim_{x \to c} f(x) = L$.
  \item For every $(x_n) \subseteq A \setminus \{c\}$ with
    $x_n \to c$, $\; f(x_n) \to L$.
\end{enumerate}

\smallskip\noindent
\hyperref[prf:sequential-criterion-limits]{\textit{Go to proof.}}
\end{proposition}

\begin{remark*}[Standard quantified statement]
\[
\begin{aligned}
(a)\quad
  &\ForallOT{\varepsilon},\;
   \ExistsIT{\delta},\;
   \forall x \in A :\;
   \operatorname{InputTolerance}(x,c,\delta)
   \Rightarrow
   \operatorname{OutputTolerance}(f(x),L,\varepsilon). \\[4pt]
(b)\quad
  &\forall (x_n) \subseteq A \setminus \{c\} :\;
   x_n \to c \;\Rightarrow\; f(x_n) \to L.
\end{aligned}
\]
\end{remark*}

\begin{remark*}[Negated quantified statement]
\[
  \neg\operatorname{FunctionLimit}(f,c,L)
  \;\iff\;
  \exists (x_n) \subseteq A \setminus \{c\} :\;
  x_n \to c \;\wedge\; f(x_n) \not\to L.
\]
\end{remark*}

\begin{remark*}[The bridge principle]
The Sequential Criterion is the engine of the algebra of limits.
Once it is in hand:
\[
  \text{algebra of sequence limits}
  \;\xRightarrow{\text{Sequential Criterion}}\;
  \text{algebra of function limits.}
\]
Every sequence result — sum, product, quotient, squeeze — transfers
to function limits by feeding a sequence into the criterion and
reading the result back. The negation is equally powerful: to
disprove $\operatorname{FunctionLimit}(f,c,L)$, exhibit one sequence
$(x_n) \subseteq A \setminus \{c\}$ with $x_n \to c$ and
$f(x_n) \not\to L$.
\end{remark*}

% ---------------------------------------------------------
% Corollary — Three Equivalent Forms
% ---------------------------------------------------------

\begin{corollary}[Three Equivalent Forms of the Function Limit]
\label{cor:limit-three-equiv}
Let $f : A \to \mathbb{R}$ and let $c$ be a cluster point of $A$.
The following are equivalent:
\begin{enumerate}[label=(\roman*)]
  \item $\operatorname{FunctionLimit}(f,c,L)$.
    \hfill ($\varepsilon$-$\delta$)
  \item $\forall V_\varepsilon(L),\;\exists V_\delta(c) :\;
    x \in V_\delta^*(c) \cap A
    \Rightarrow f(x) \in V_\varepsilon(L)$.
    \hfill (neighbourhood)
  \item $\forall (x_n) \subseteq A \setminus \{c\},\;
    x_n \to c \Rightarrow f(x_n) \to L$.
    \hfill (sequential)
\end{enumerate}
\end{corollary}

\begin{remark*}[Interpretation]
All three encodings capture the same mathematical condition. The
$\varepsilon$-$\delta$ form is best for proving explicit bounds. The
neighbourhood form is best for geometric reasoning and metric space
generalisation. The sequential form is best for importing sequence
results and constructing counterexamples.
\end{remark*}

% ---------------------------------------------------------
% Theorem — Limit Implies Local Boundedness
% ---------------------------------------------------------

\begin{theorem}[Limit Implies Local Boundedness]
\label{thm:limit-implies-local-bounding}
Let $f : A \to \mathbb{R}$ and let $c$ be a cluster point of $A$.
If $\operatorname{FunctionLimit}(f, c, L)$ for some $L \in \mathbb{R}$,
then $\operatorname{BoundedNear}(f, c)$: there exist $\delta, M > 0$
such that $|f(x)| \leq M$ for all $x \in A$ with
$\operatorname{InputTolerance}(x,c,\delta)$.

\smallskip\noindent
\hyperref[prf:limit-implies-local-bounding]{\textit{Go to proof.}}
\end{theorem}

\begin{remark*}[Standard quantified statement]
\[
  \operatorname{FunctionLimit}(f,c,L)
  \;\Rightarrow\;
  \operatorname{BoundedNear}(f,c)
  \;\colonequiv\;
  \exists\delta > 0,\;\exists M > 0,\;
  \forall x \in A :\;
  0 < |x-c| < \delta \Rightarrow |f(x)| \leq M.
\]
\end{remark*}

\begin{remark*}[Proof sketch]
Apply the definition with $\varepsilon = 1$: get $\delta > 0$ such
that $\operatorname{OutputTolerance}(f(x),L,1)$ holds for all $x$
satisfying $\operatorname{InputTolerance}(x,c,\delta)$. The triangle
inequality then gives $|f(x)| \leq |f(x)-L| + |L| < 1 + |L|$.
Set $M = 1 + |L|$.
\end{remark*}

\begin{remark*}[Interpretation]
A finite limit does more than control the output near $L$: it
prevents the function from growing without bound near $c$. Local
boundedness is a coarser statement than limit existence, but it is
what is needed for the product and quotient rules.
\end{remark*}

% ---------------------------------------------------------
% Theorem — Bounded Limits
% ---------------------------------------------------------

\begin{theorem}[Bounded Limits]
\label{thm:bounded-limits}
Let $f : A \to \mathbb{R}$, let $c$ be a cluster point of $A$, and
let $a, b \in \mathbb{R}$. If $a \leq f(x) \leq b$ for all $x \in A$
with $x \neq c$, and $\operatorname{FunctionLimit}(f,c,L)$, then
$a \leq L \leq b$.

\smallskip\noindent
\hyperref[prf:bounded-limits]{\textit{Go to proof.}}
\end{theorem}

\begin{remark*}[Standard quantified statement]
\[
  \bigl(\forall x \in A \setminus \{c\} :\; a \leq f(x) \leq b\bigr)
  \;\wedge\;
  \operatorname{FunctionLimit}(f,c,L)
  \;\Rightarrow\;
  a \leq L \leq b.
\]
\end{remark*}

\begin{remark*}[Interpretation]
If every nearby value of $f$ is trapped in $[a,b]$, the limit cannot
escape outside $[a,b]$. This is the function-limit analogue of the
sequence result that limits preserve non-strict inequalities. The
proof applies the Sequential Criterion and that sequence result
directly.
\end{remark*}

% ---------------------------------------------------------
% Theorem — Squeeze Theorem for Function Limits
% ---------------------------------------------------------

\begin{theorem}[Squeeze Theorem for Function Limits]
\label{thm:squeeze-function-limits}
Let $f, g, h : A \to \mathbb{R}$ and let $c$ be a cluster point of
$A$. Suppose $\operatorname{FunctionLimit}(f,c,L) =
\operatorname{FunctionLimit}(g,c,L)$, and suppose there exists
$\delta_0 > 0$ such that
\[
  \forall x \in A,\;
  0 < |x - c| < \delta_0
  \;\Rightarrow\;
  f(x) \leq h(x) \leq g(x).
\]
Then $\operatorname{FunctionLimit}(h, c, L)$.

\smallskip\noindent
\hyperref[prf:squeeze-function-limits]{\textit{Go to proof.}}
\end{theorem}

\begin{remark*}[Standard quantified statement]
\[
  \operatorname{FunctionLimit}(f,c,L)
  = \operatorname{FunctionLimit}(g,c,L)
  \;\wedge\;
  f \leq h \leq g \text{ near } c
  \;\Rightarrow\;
  \operatorname{FunctionLimit}(h,c,L).
\]
\end{remark*}

\begin{remark*}[Interpretation]
The proof routes entirely through the Sequential Criterion: convert
to sequences, apply the sequential squeeze theorem, convert back. No
new $\varepsilon$-$\delta$ argument is needed. The condition
$f \leq h \leq g$ is only required near $c$, not globally.
\end{remark*}

% ---------------------------------------------------------
% Theorem — Limits Are Local
% ---------------------------------------------------------

\begin{theorem}[Limits Are Local]
\label{thm:limits-are-local}
Let $E \subseteq X \subseteq \mathbb{R}$, let $x_0$ be a cluster
point of $E$, let $f : X \to \mathbb{R}$, and let $\delta > 0$.
Then
\[
  \operatorname{FunctionLimit}(f\!\restriction_E,\, x_0,\, L)
  \;\iff\;
  \operatorname{FunctionLimit}(f\!\restriction_{E \cap (x_0-\delta,\,
  x_0+\delta)},\, x_0,\, L).
\]

\smallskip\noindent
\hyperref[prf:limits-are-local]{\textit{Go to proof.}}
\end{theorem}

\begin{remark*}[Standard quantified statement]
\[
  \lim_{x \to x_0;\, x \in E} f(x) = L
  \;\iff\;
  \lim_{x \to x_0;\, x \in E \cap (x_0 - \delta,\, x_0 + \delta)}
  f(x) = L.
\]
\end{remark*}

\begin{remark*}[Interpretation]
The existence and value of $\operatorname{FunctionLimit}(f, x_0, L)$
depend only on the behavior of $f$ within any fixed neighbourhood of
$x_0$. Points outside that neighbourhood are irrelevant. This
justifies focusing all limit arguments on small balls around $c$
without any loss of generality.
\end{remark*}

% ---------------------------------------------------------
% Proposition — Limit of Absolute Value
% ---------------------------------------------------------

\begin{proposition}[Limit of Absolute Value]
\label{prop:limit-absolute-value}
Let $f : A \to \mathbb{R}$ and let $c$ be a cluster point of $A$.
If $\operatorname{FunctionLimit}(f,c,L)$, then
$\operatorname{FunctionLimit}(|f|, c, |L|)$.

\smallskip\noindent
\hyperref[prf:limit-absolute-value]{\textit{Go to proof.}}
\end{proposition}

\begin{remark*}[Standard quantified statement]
\[
  \operatorname{FunctionLimit}(f,c,L)
  \;\Rightarrow\;
  \operatorname{FunctionLimit}(|f|, c, |L|).
\]
\end{remark*}

\begin{remark*}[Negated quantified statement]
The converse fails. $\operatorname{FunctionLimit}(|f|,c,|L|)$ does
not imply $\operatorname{FunctionLimit}(f,c,L)$. Counterexample:
$f(x) = \operatorname{sgn}(x)$ at $c = 0$ has $|f(x)| \equiv 1$ so
$\lim|f| = 1$, but $\lim f$ does not exist.
\end{remark*}

\begin{remark*}[Proof strategy]
The reverse triangle inequality
$\bigl||f(x)| - |L|\bigr| \leq |f(x) - L|$
feeds the $\varepsilon$-$\delta$ definition directly: the same
$\delta$ that delivers $\operatorname{OutputTolerance}(f(x),L,\varepsilon)$
also delivers $\operatorname{OutputTolerance}(|f(x)|,|L|,\varepsilon)$.
\end{remark*}

% ---------------------------------------------------------
% Theorem — Composition of Limits
% ---------------------------------------------------------

\begin{theorem}[Composition of Limits]
\label{thm:composition-of-limits}
Let $f : A \to B$, let $g : B \to \mathbb{R}$, let $c_1$ be a
cluster point of $A$, and let $c_2 \in \mathbb{R}$ be a cluster
point of $B$. Suppose $\operatorname{FunctionLimit}(f, c_1, c_2)$
and $\operatorname{FunctionLimit}(g, c_2, L)$. If $c_2 \in B$,
assume also that $g(c_2) = L$. Then
$\operatorname{FunctionLimit}(g \circ f, c_1, L)$.

\smallskip\noindent
\hyperref[prf:composition-of-limits]{\textit{Go to proof.}}
\end{theorem}

\begin{remark*}[Standard quantified statement]
\[
  \operatorname{FunctionLimit}(f,c_1,c_2)
  \;\wedge\;
  \operatorname{FunctionLimit}(g,c_2,L)
  \;\wedge\;
  \bigl(c_2 \in B \Rightarrow g(c_2) = L\bigr)
  \;\Rightarrow\;
  \operatorname{FunctionLimit}(g \circ f,\, c_1,\, L).
\]
\end{remark*}

\begin{remark*}[Failure mode decomposition]
The hypothesis $g(c_2) = L$ is essential. Without it the composition
can fail: take $f \equiv c_2$ (constant) and $g(c_2) \neq L$.
Then $g(f(x)) = g(c_2) \neq L$ for all $x$, so
$\operatorname{FunctionLimit}(g \circ f, c_1, L)$ fails even
though both component limits exist. The failure is that $f(x) = c_2$
is excluded from the punctured condition on $g$, but $f$ hits $c_2$
constantly.
\end{remark*}

\begin{remark*}[Interpretation]
Composition of limits requires care precisely because the input
condition is punctured. The extra hypothesis $g(c_2) = L$ bridges the
gap: when $f(x) = c_2$, the output condition $g(f(x)) = g(c_2) = L$
is satisfied directly. This hypothesis is equivalent to requiring $g$
to be continuous at $c_2$ (within $B$).
\end{remark*}

% ---------------------------------------------------------
% Theorem — Monotone Limit Theorem
% ---------------------------------------------------------

\begin{theorem}[Monotone Limit Theorem]
\label{thm:monotone-limit-theorem}
Let $a < b$ and let $f : (a,b) \to \mathbb{R}$ be monotone and
bounded. Then both one-sided limits exist at every interior point.
For increasing $f$ and every $x_0 \in (a,b)$:
\[
  f(x_0^-) = \sup\{\,f(x) : a < x < x_0\,\},
  \qquad
  f(x_0^+) = \inf\{\,f(x) : x_0 < x < b\,\}.
\]

\smallskip\noindent
\hyperref[prf:monotone-limit-theorem]{\textit{Go to proof.}}
\end{theorem}

\begin{remark*}[Standard quantified statement]
\[
  \operatorname{MonotoneFunction}(f,(a,b))
  \;\wedge\;
  \operatorname{BoundedOn}(f,(a,b))
  \;\Rightarrow\;
  \forall x_0 \in (a,b) :\;
  \operatorname{LeftLimit}(f,x_0,f(x_0^-))
  \;\wedge\;
  \operatorname{RightLimit}(f,x_0,f(x_0^+)).
\]
\end{remark*}

\begin{remark*}[Interpretation]
Boundedness guarantees the supremum and infimum used as candidate
limits exist. Monotonicity forces convergence toward them from each
side: for increasing $f$, the values $f(x)$ for $x < x_0$ form a
bounded increasing net approaching their supremum. No oscillation is
possible. The theorem establishes that monotone functions can only
have jump discontinuities — never oscillatory ones.
\end{remark*}

% ---------------------------------------------------------
% Proposition — Cauchy Criterion for Function Limits
% ---------------------------------------------------------

\begin{proposition}[Cauchy Criterion for Function Limits]
\label{prop:cauchy-criterion-limits}
Let $f : A \to \mathbb{R}$ and let $c$ be a cluster point of $A$.
Then $\operatorname{FunctionLimit}(f,c,L)$ exists for some
$L \in \mathbb{R}$ if and only if
\[
  \ForallOT{\varepsilon},\;
  \ExistsIT{\delta},\;
  \forall x, y \in A :\;
  \operatorname{InputTolerance}(x,c,\delta)
  \;\wedge\;
  \operatorname{InputTolerance}(y,c,\delta)
  \;\Rightarrow\;
  \operatorname{OutputTolerance}(f(x),f(y),\varepsilon).
\]

\smallskip\noindent
\hyperref[prf:cauchy-criterion-limits]{\textit{Go to proof.}}
\end{proposition}

\begin{remark*}[Standard quantified statement]
\[
  \exists L \in \mathbb{R},\;
  \operatorname{FunctionLimit}(f,c,L)
  \;\iff\;
  \forall\varepsilon>0,\;\exists\delta>0,\;
  \forall x,y \in A :\;
  0<|x-c|<\delta \wedge 0<|y-c|<\delta
  \Rightarrow |f(x)-f(y)|<\varepsilon.
\]
\end{remark*}

\begin{remark*}[Interpretation]
The Cauchy criterion tests limit existence without knowing $L$. Two
outputs from any sufficiently small punctured ball must lie within
$\varepsilon$ of each other — the outputs cluster even if their
common limit is unknown. This is the function-limit analogue of the
Cauchy completeness criterion for sequences, and relies on the
completeness of $\mathbb{R}$ in exactly the same way.
\end{remark*}
